% META: TODO.1

\section{Экспериментальная проверка и анализ результатов}

\subsection{План экспериментов}
Для оценки качества разработанного подхода были проведены следующие эксперименты:

\begin{enumerate}
    \item \textbf{Обучение классификаторов} для каждой из четырёх операций на соответствующих подвыборках. Оценка - accuracy на тестовой выборке (10\% от датасета).
    \item \textbf{Обучение сегментаторов} для каждой операции. Оценка - доля верно классифицированных рёбер (edge accuracy, для каждого класса), доля моделей с accuracy >= 90\%, >=50\%, >=10\%, и метрика IoU для каждого класса на тестовой выборке (10\% от датасета).
    \item \textbf{Сравнение с baseline}: тривиальный классификатор, всегда предсказывающий доминирующий класс.
    \item \textbf{Тестирование на реальных деталях} из (TODO: proper appendix name)Приложения Б: визуальная оценка качества сегментации и успешность упрощения в КОМПАС-3D.
\end{enumerate}

\subsection{Результаты обучения классификаторов} (TODO: whole subsection below, proper figs and tables)
Графики обучения (рис.~\ref{fig:acc-loss}) показывают устойчивую сходимость: точность на тесте растёт, а функция потерь на обучении монотонно убывает. 

\begin{figure}[h]
    \centering
    % TODO: Вставить графики из отчёта: Рисунок 4а и 4б
    \includegraphics[width=0.45\linewidth]{TODO_plot_acc.png}
    \includegraphics[width=0.45\linewidth]{TODO_plot_loss.png}
    \caption{Точность на тестовой выборке (слева) и функция потерь на обучении (справа) для операций bossExtrusion и fillet}
    \label{fig:acc-loss}
\end{figure}

Численные результаты классификации сведены в таблицу~\ref{tab:cls-results}.

\begin{table}[h]
\centering
\caption{Результаты классификации (accuracy на тесте)}
\label{tab:cls-results}
\begin{tabular}{|l|c|c|}
\hline
\textbf{Операция} & \textbf{Accuracy (наша модель)} & \textbf{Baseline (доминирующий класс)} \\
\hline
bossExtrusion & 84.0\% & 50\% (сбалансированная выборка) \\
\hline
cutExtrusion & 67.9\% & 50\% \\
\hline
fillet & 90.6\% & 50\% \\
\hline
chamfer & 70.1\% & 50\% \\
\hline
\end{tabular}
\end{table}

Наилучший результат достигнут для операции \texttt{fillet} (90.6\%), что объясняется характерной геометрией скруглений. Наихудший - для \texttt{cutExtrusion} (67.9\%), поскольку вырезания часто маскируются последующими операциями.

\subsection{Результаты сегментации} (TODO: whole subsection below, proper figs and tables)
Точность сегментации (доля верно классифицированных рёбер) и доля моделей с качеством >= 90\% приведены в таблице~\ref{tab:seg-results}.

\begin{table}[h]
\centering
\caption{Результаты сегментации}
\label{tab:seg-results}
\begin{tabular}{|l|c|c|c|}
\hline
\textbf{Операция} & \textbf{Edge accuracy} & \textbf{Доля моделей с acc >= 90\%} & \textbf{Размер выборки (тест)} \\
\hline
bossExtrusion & 78.4\% & 46.8\% & TODO \\
\hline
cutExtrusion (curve) & 96.8\% & 91.3\% & TODO \\
\hline
fillet & 93.6\% & 78.7\% & TODO \\
\hline
chamfer & 68.1\% & 53.1\% & TODO \\
\hline
\end{tabular}
\end{table}

Операция \texttt{cutExtrusion} на криволинейных контурах демонстрирует выдающуюся точность (96.8\%), вероятно, из-за чётких границ вырезания. \texttt{chamfer} показывает худший результат из-за малого числа примеров и схожести с фасками других размеров.

\subsection{Анализ типовых ошибок} (TODO: whole subsection below, proper figs and tables)
Визуальный анализ сегментации (рис.~\ref{fig:seg-errors}) выявил следующие характерные ошибки:

\begin{figure}[h]
    \centering
    % TODO: Вставить примеры ошибок из отчёта (отверстия как скругления, ложноположительные рёбра)
    \includegraphics[width=0.3\linewidth]{TODO_error_hole_as_fillet.png}
    \includegraphics[width=0.3\linewidth]{TODO_error_false_positive.png}
    \caption{Типовые ошибки: отверстие, классифицированное как скругление (слева); ложноположительные рёбра на плоской грани (справа)}
    \label{fig:seg-errors}
\end{figure}

\begin{itemize}
    \item \textbf{Отверстия, ошибочно принимаемые за скругления}: геометрические признаки цилиндрических поверхностей и скруглений частично пересекаются, особенно при низком разрешении сетки.
    \item \textbf{Ложноположительные рёбра на плоских гранях}: возникают из-за шумов триангуляции или несовершенства разметки.
    \item \textbf{Пропуск мелких операций}: скругления малого радиуса или фаски на сложных поверхностях часто не обнаруживаются.
\end{itemize}

\subsection{Тестирование на реальных деталях} (TODO: whole subsection below, proper figs and tables)
Модуль \texttt{main.exe} был протестирован на пяти типовых деталях (Приложение Б). Во всех случаях удалось выделить геометрию, соответствующую заданной операции, с визуально приемлемой точностью. Пример для детали «Деталь 01» показан на рисунке~\ref{fig:real-part}.

\begin{figure}[h]
    \centering
    % TODO: Вставить скриншоты из Приложения Б (исходная модель, обнаруженная часть, оставшаяся часть)
    \includegraphics[width=0.3\linewidth]{TODO_part01_original.png}
    \includegraphics[width=0.3\linewidth]{TODO_part01_detected.png}
    \includegraphics[width=0.3\linewidth]{TODO_part01_remaining.png}
    \caption{Пример работы модуля на детали «Деталь 01»: исходная модель, обнаруженная часть (fillet), оставшаяся геометрия}
    \label{fig:real-part}
\end{figure}

Интеграция с КОМПАС-3D позволила успешно упростить модели в автоматическом режиме для 8 из 10 тестовых случаев. Неудачи были связаны с чрезмерно сложной геометрией (>2000 треугольников) или отсутствием обученной операции (например, вращение).










\subsection{Обучение независимых моделей по операциям}
Для каждой из четырёх операций (\texttt{bossExtrusion}, \texttt{cutExtrusion}, \texttt{fillet}, \texttt{chamfer}) были обучены отдельные классификатор и сегментатор. Такой подход позволяет:
\begin{itemize}
    \item использовать специфичные для операции гиперпараметры (что, как будет показано далее, вполне практическое требование);
    \item дообучать или заменять модель для одной операции, не затрагивая другие;
    \item легко добавлять новые операции в будущем.
\end{itemize}

\subsubsection{Проведённые серии экспериментов}
В ходе работы было выполнено несколько запусков для классификации и сегментации с различными архитектурными конфигурациями. Ниже перечислены основные комбинации гиперпараметров:

\begin{itemize}
    \item \textbf{Классификация, конфигурация A}: 
        \begin{itemize}
            \item \texttt{--ncf 32 32 64 64 128}
            \item \texttt{--pool\_res 2400 2200 2000 1800 1600}
            \item \texttt{--resblocks 3}
        \end{itemize}
        Результаты:
        \begin{itemize}
            \item \texttt{}
            \item \texttt{}
            \item \texttt{}
        \end{itemize}
    \item \textbf{Классификация, конфигурация B}: 
        \begin{itemize}
            \item \texttt{--ncf 32 64 128}
            \item \texttt{--pool\_res 2400 1800 1200}
            \item \texttt{--resblocks 1}
        \end{itemize}
    \item \textbf{Сегментация, конфигурация C}: 
        \begin{itemize}
            \item \texttt{--ncf 32 32 64 64 128 128}
            \item \texttt{--pool\_res 2400 2200 2000 1800 1600}
            \item \texttt{--resblocks 3}
        \end{itemize}
        Результаты:
          cutExtrusion:
            \begin{itemize}
                \item \texttt{Accuracy (общая): 0.854}
                \item \texttt{Accuracy (выбранная операция): 0.758}
                \item \texttt{Accuracy (остальные операции): 0.883}
                \item \texttt{IoU (выбранная операция): 0.547}
                \item \texttt{IoU (остальные операции): 0.820}
                \item \texttt{Amount where (accuracy >= 10\%): 1.0}
                \item \texttt{Amount where (accuracy >= 50\%): 0.957}
                \item \texttt{Amount where (accuracy >= 90\%): 0.551}
            \end{itemize}
          fillet:
            \begin{itemize}
                \item \texttt{Accuracy (общая): 0.895}
                \item \texttt{Accuracy (выбранная операция): 0.917}
                \item \texttt{Accuracy (остальные операции): 0.886}
                \item \texttt{IoU (выбранная операция): 0.716}
                \item \texttt{IoU (остальные операции): 0.856}
                \item \texttt{Amount where (accuracy >= 10\%): 1.0}
                \item \texttt{Amount where (accuracy >= 50\%): 1.0}
                \item \texttt{Amount where (accuracy >= 90\%): 0.609}
            \end{itemize}
          chamfer:
            \begin{itemize}
                \item \texttt{Accuracy (общая): 0.959}
                \item \texttt{Accuracy (выбранная операция): 0.941}
                \item \texttt{Accuracy (остальные операции): 0.964}
                \item \texttt{IoU (выбранная операция): 0.824}
                \item \texttt{IoU (остальные операции): 0.946}
                \item \texttt{Amount where (accuracy >= 10\%): 1.0}
                \item \texttt{Amount where (accuracy >= 50\%): 0.993}
                \item \texttt{Amount where (accuracy >= 90\%): 0.869}
            \end{itemize}
          bossExtrusion:
            \begin{itemize}
                \item \texttt{Accuracy (общая): 0.724}
                \item \texttt{Accuracy (выбранная операция): 0.679}
                \item \texttt{Accuracy (остальные операции): 0.743}
                \item \texttt{IoU (выбранная операция): 0.416}
                \item \texttt{IoU (остальные операции): 0.659}
                \item \texttt{Amount where (accuracy >= 10\%): 1.0}
                \item \texttt{Amount where (accuracy >= 50\%): 0.880}
                \item \texttt{Amount where (accuracy >= 90\%): 0.150}
            \end{itemize}
    \item \textbf{Сегментация, конфигурация D}: 
        \begin{itemize}
            \item \texttt{--ncf 32 32 64 64 128}
            \item \texttt{--pool\_res 2400 1800 1200 600}
            \item \texttt{--resblocks 1}
        \end{itemize}
        Результаты:
          cutExtrusion:
            \begin{itemize}
                \item \texttt{Accuracy (общая): 0.852}
                \item \texttt{Accuracy (выбранная операция): 0.769}
                \item \texttt{Accuracy (остальные операции): 0.878}
                \item \texttt{IoU (выбранная операция): 0.542}
                \item \texttt{IoU (остальные операции): 0.818}
                \item \texttt{Amount where (accuracy >= 10\%): 1.0}
                \item \texttt{Amount where (accuracy >= 50\%): 0.964}
                \item \texttt{Amount where (accuracy >= 90\%): 0.528}
            \end{itemize}
          fillet:
            \begin{itemize}
                \item \texttt{Accuracy (общая): 0.934}
                \item \texttt{Accuracy (выбранная операция): 0.929}
                \item \texttt{Accuracy (остальные операции): 0.937}
                \item \texttt{IoU (выбранная операция): 0.807}
                \item \texttt{IoU (остальные операции): 0.904}
                \item \texttt{Amount where (accuracy >= 10\%): 1.0}
                \item \texttt{Amount where (accuracy >= 50\%): 1.0}
                \item \texttt{Amount where (accuracy >= 90\%): 0.794}
            \end{itemize}
          chamfer:
            \begin{itemize}
                \item \texttt{Accuracy (общая): 0.962}
                \item \texttt{Accuracy (выбранная операция): 0.929}
                \item \texttt{Accuracy (остальные операции): 0.970}
                \item \texttt{IoU (выбранная операция): 0.840}
                \item \texttt{IoU (остальные операции): 0.948}
                \item \texttt{Amount where (accuracy >= 10\%): 1.0}
                \item \texttt{Amount where (accuracy >= 50\%): 0.993}
                \item \texttt{Amount where (accuracy >= 90\%): 0.883}
            \end{itemize}
          bossExtrusion:
            \begin{itemize}
                \item \texttt{Accuracy (общая): 0.744}
                \item \texttt{Accuracy (выбранная операция): 0.748}
                \item \texttt{Accuracy (остальные операции): 0.742}
                \item \texttt{IoU (выбранная операция): 0.458}
                \item \texttt{IoU (остальные операции): 0.676}
                \item \texttt{Amount where (accuracy >= 10\%): 1.0}
                \item \texttt{Amount where (accuracy >= 50\%): 0.910}
                \item \texttt{Amount where (accuracy >= 90\%): 0.248}
            \end{itemize}
    \item \textbf{Сегментация, конфигурация E}: 
        \begin{itemize}
            \item \texttt{--ncf 32 64 64 128}
            \item \texttt{--pool\_res 2400 1800 1200}
            \item \texttt{--resblocks 1}
        \end{itemize}
        Результаты:
          cutExtrusion:
            \begin{itemize}
                \item \texttt{Accuracy (общая): 0.851}
                \item \texttt{Accuracy (выбранная операция): 0.749}
                \item \texttt{Accuracy (остальные операции): 0.882}
                \item \texttt{IoU (выбранная операция): 0.533}
                \item \texttt{IoU (остальные операции): 0.817}
                \item \texttt{Amount where (accuracy >= 10\%): 1.0}
                \item \texttt{Amount where (accuracy >= 50\%): 0.964}
                \item \texttt{Amount where (accuracy >= 90\%): 0.528}
            \end{itemize}
          fillet:
            \begin{itemize}
                \item \texttt{Accuracy (общая): 0.941}
                \item \texttt{Accuracy (выбранная операция): 0.934}
                \item \texttt{Accuracy (остальные операции): 0.944}
                \item \texttt{IoU (выбранная операция): 0.830}
                \item \texttt{IoU (остальные операции): 0.918}
                \item \texttt{Amount where (accuracy >= 10\%): 1.0}
                \item \texttt{Amount where (accuracy >= 50\%): 1.0}
                \item \texttt{Amount where (accuracy >= 90\%): 0.837}
            \end{itemize}
          chamfer:
            \begin{itemize}
                \item \texttt{Accuracy (общая): 0.963}
                \item \texttt{Accuracy (выбранная операция): 0.942}
                \item \texttt{Accuracy (остальные операции): 0.968}
                \item \texttt{IoU (выбранная операция): 0.836}
                \item \texttt{IoU (остальные операции): 0.949}
                \item \texttt{Amount where (accuracy >= 10\%): 1.0}
                \item \texttt{Amount where (accuracy >= 50\%): 0.993}
                \item \texttt{Amount where (accuracy >= 90\%): 0.883}
            \end{itemize}
          bossExtrusion:
            \begin{itemize}
                \item \texttt{Accuracy (общая): 0.710}
                \item \texttt{Accuracy (выбранная операция): 0.739}
                \item \texttt{Accuracy (остальные операции): 0.699}
                \item \texttt{IoU (выбранная операция): 0.423}
                \item \texttt{IoU (остальные операции): 0.635}
                \item \texttt{Amount where (accuracy >= 10\%): 1.0}
                \item \texttt{Amount where (accuracy >= 50\%): 0.827}
                \item \texttt{Amount where (accuracy >= 90\%): 0.165}
            \end{itemize}
\end{itemize}

Каждая конфигурация обучалась на одних и тех же наборах данных (по операциям) не менее трёх раз со случайными начальными весами. В качестве итоговой бралась модель, показавшая наилучшую точность на валидационной выборке. Сравнение результатов по этим конфигурациям позволило выбрать оптимальную архитектуру для каждой операции: ....

\subsection{Обучение независимых моделей по операциям}
Для каждой из четырёх операций (\texttt{bossExtrusion}, \texttt{cutExtrusion}, \texttt{fillet}, \texttt{chamfer}) были обучены отдельные классификатор и сегментатор. Такой подход позволяет:
\begin{itemize}
    \item использовать специфичные для операции гиперпараметры;
    \item дообучать или заменять модель для одной операции, не затрагивая другие;
    \item легко добавлять новые операции в будущем.
\end{itemize}

